\subsection{ALTERNATIVE ALGEBRAS}

Let A be a commutative ring and F a (not necessarily associative) A-algebra. For any three elements x, y, z of F, we write

\[
a (x, y, z) = x (y z) - (x y) z\tag{1}
\]

(associator of \(x, y, z\) ); \(a\) is obviously an A-trilinear mapping of \(\mathbf{F} \times \mathbf{F} \times \mathbf{F}\) into \(\mathbf{F}\) .

Lemma 1. For all \(p, q, r, s\) in the algebra \(\mathbf{F}\) ,

\[
a (p q, r, s) - a (p, q r, s) + a (p, q, r s) = a (p, q, r) s + p a (q, r, s).\tag{2}
\]

The verification follows immediately from definition (1).

PROPOSITION 1. For an A-algebra F, the following conditions are equivalent:

\begin{enumerate}
\def\labelenumi{(\alph{enumi})}
\item
  For every ordered pair of elements \(x, y\) of \(\mathbf{F}\) , the subalgebra generated by \(x\) and \(y\) is associative.
\item
  The trilinear mapping \((x, y, z) \mapsto a(x, y, z)\) is alternating (§ 7, no. 3).
\item
  For every ordered pair of elements \(x, y\) of \(\mathbf{F}\) , \(x^2y = x(xy)\) and \(yx^2 = (yx)x\) .
\end{enumerate}

Clearly (a) implies (c). We show that (c) implies (b): by definition (§ 7, no. 3) to prove (b), it suffices to verify that \(a(x, x, y) = 0\) and \(a(x, y, y) = 0\) , which is precisely (c).

To prove that (b) implies (a), we use the following 4 lemmas:

Lemma 2. Let E be an A-algebra such that the trilinear mapping \((x,y,z)\mapsto a(x,y,z)\) is alternating, S a generating system of E and U a sub-A-module of E containing S and such that \(sU\subset U\) and \(Us\subset U\) for all \(s\in S\) . Then U = E.

The set \(U'\) of \(x \in E\) such that \(xU \subset U\) and \(Ux \subset U\) is obviously a sub-A-module of E, which contains S by hypothesis. On the other hand, for x, y in \(U'\) and \(u \in U\) , by hypothesis

\[
(x y) u = x (y u) + a (x, y, u) = x (y u) - a (x, u, y) = x (y u) - (x u) y + x (u y) \in \mathrm{U};
\]

on passing to the opposite algebra, we have similarly \(u(xy) \in \mathbf{U}\) . Hence \(\mathbf{U}'\) is a subalgebra of \(\mathbf{E}\) and, since it contains \(\mathbf{S}\) , \(\mathbf{U}' = \mathbf{E}\) . Hence \(\mathbf{EU} \subset \mathbf{U}\) and a fortiori \(\mathbf{UU} \subset \mathbf{U}\) , which proves that \(\mathbf{U}\) is a subalgebra of \(\mathbf{E}\) ; as it contains \(\mathbf{S}\) , \(\mathbf{U} = \mathbf{E}\) , which proves the lemma.

A subset H of F is called strongly associative if \(a(u, v, w) = 0\) when at least two of the elements u, v, w belong to H.

Lemma 3. Suppose that the mapping \(a\) is alternating. If \(\mathbf{H}\) is a strongly associative subset of \(\mathbf{F}\) , the subalgebra of \(\mathbf{F}\) generated by \(\mathbf{H}\) is strongly associative.

As the set of strongly associative subsets of F is inductive, it suffices to prove that if H is a maximal strongly associative subset of F, H is then a subalgebra of F. As H is obviously a sub-A-module of F, it suffices to verify that for any two elements u, v of H, \(H \cup \{uv\}\) is also strongly associative, for by virtue of the definition of H, this will imply \(uv \in H\) . Now, for all \(z \in H\) and all \(t \in F\) , by (2)

\[
a (u v, t, z) - a (u, v t, z) + a (u, v, t z) = 0
\]

since \(\mathbf{H}\) is strongly associative; as \(u, v, z\) are in \(\mathbf{H}\) , also

\[
a (u, v t, z) = a (u, v, t z) = 0,
\]

whence \(a(uv, t, z) = 0\) . Using the fact that a is alternating, this shows that \(a(p, q, r) = 0\) whenever at least two of the elements p, q, r belong to \(H \cup \{uv\}\) whence the lemma.

Lemma 4. Suppose that the mapping \(a\) is alternating. Then, for all \(x \in \mathbf{F}\) , the subalgebra of \(\mathbf{F}\) generated by \(x\) is strongly associative.

\(a(u,v,w) = 0\) whenever two of the three elements \(u, v, w\) are equal to \(x\) and it suffices to apply Lemma 3.

Lemma 5. Suppose that the mapping \(a\) is alternating and let \(\mathbf{X}, \mathbf{Y}\) be two strongly associative subalgebras of \(\mathbf{F}\) . Then the subalgebra of \(\mathbf{E}\) generated by \(\mathbf{X} \cup \mathbf{Y}\) is associative.

Let Z be the set of \(z \in E\) such that \(a(u, v, z) = 0\) for all \(u \in X\) and \(v \in Y\) , this is obviously a sub-A-module containing X and Y since X and Y are strongly associative; by Lemma 2, it will suffice to verify that, for \(u \in X\) and \(v \in Y\) , \(uZ \subset Z\) , \(vZ \subset Z\) , \(Zu \subset Z\) and \(Zv \subset Z\) . Now, for \(u, u'\) in X, \(v \in Y\) and \(z \in Z\) , by (2)

\[
a (u ^ {\prime} u, z, v) - a (u ^ {\prime}, u z, v) + a (u ^ {\prime}, u, z v) = a (u ^ {\prime}, u, z) v + u ^ {\prime} a (u, z, v) = 0
\]

by virtue of the fact that X is strongly associative and the definition of Z. But as X is strongly associative, \(a(u', u, zv) = 0\) and since \(u'u \in X\) , \(a(u'u, z, v) = 0\) by definition of Z. Hence \(a(u', uz, v) = 0\) , which shows that \(uZ \subset Z\) . Applying (2) now with \((p, q, r, s) = (v, z, u, u')\) , we obtain similarly \(Zu \subset Z\) . Interchanging the roles of X and Y and using the fact that a is alternating, we obtain \(vZ \subset Z\) and \(Zv \subset Z\) ; whence the lemma.

It now suffices, to prove that (b) implies (a) in Proposition 1, to take \(\mathbf{X} = \{x\}\) and \(\mathbf{Y} = \{y\}\) , using Lemma 4.

DEFINITION 1. An algebra \(\mathbf{F}\) is called alternative if it satisfies the equivalent conditions of Proposition 1.

An associative algebra is obviously alternative. In no. 3 we shall give an example of an alternative algebra which is not associative.

If \(\mathbf{F}\) is an alternative A-algebra, every A-algebra \(\mathbf{F} \otimes_{\mathbf{A}} \mathbf{A}'\) obtained from \(\mathbf{F}\) by extending the scalars (§ 1, no. 5) is an alternative \(\mathbf{A}'\) -algebra, as follows from condition (b) of Proposition 1.

